\chapter{The Integral Model and the Model Seam}
\label{chap:lwx-model}

Theorem~\ref{thm:lwx316} is an estimate about a matrix.  For it to say anything about modular
forms one has to know two things: that the matrix really is the matrix of $U_p$ on a space of
overconvergent automorphic forms, and that its Fredholm determinant really is the Fredholm
determinant of that operator.  These are \cite[Proposition 3.1]{LWX} and
\cite[Proposition 2.17]{LWX}, and they are what this chapter proves.

The difficulty is that the two sides use different coordinates.  The integral side of
Chapter~\ref{chap:lwx-halo} works in the Mahler basis $\binom{z}{n}$ of $C(\Zp, \Lhalo)$, in
which the entries are integral but the operator is not visibly compact; the analytic side of
Chapter~\ref{chap:qmf} works in the monomial basis of the Tate algebra, in which the operator is
visibly compact but the entries are not integral.  \cite[proof of Proposition 2.17]{LWX} bridges
them in one sentence: ``$P$ and $P'$ are conjugated by an infinite diagonal matrix with entries
$\lfloor n/(q^{-1}p^{m})\rfloor!$''.  The intermediate basis is Colmez's, and the statement that
it is orthonormal is Amice's theorem.  So the chapter is: analysis
(\S\ref{sec:explog}--\S\ref{sec:haloweight}), the integral model (\S\ref{sec:intmodel}), Amice's
theorem (\S\ref{sec:amice}), and the seam (\S\ref{sec:seam}).

\section{$p$-adic analysis over an abstract field}
\label{sec:explog}

\begin{definition}
  \label{def:padicexplog}
  \lean{LWX.PadicExpLog.padicLog, LWX.PadicExpLog.padicExp, LWX.PadicExpLog.padicLog_mul,
        LWX.PadicExpLog.padicExp_add, LWX.PadicExpLog.padicExp_padicLog,
        LWX.PadicExpLog.padicLog_padicExp}
  \leanok
  Over a complete ultrametric field $K$ with $\lVert p\rVert<1$ and $p$ odd,
  $\log u = \sum_{n\ge0}(-1)^{n}\tfrac{(u-1)^{n+1}}{n+1}$ for $\lVert u-1\rVert<1$ and
  $\exp w = \sum_n w^{n}/n!$ for $\lVert w\rVert^{2}<\lVert p\rVert$.  On the joint disc they are
  mutually inverse, and $\log$ is additive.
\end{definition}

The convergence discs are \cite[Chapter IV \S2]{Koblitz}.  Additivity of $\log$ is obtained from
$\exp\circ\log = \mathrm{id}$ and injectivity rather than by rearranging series.  Odd $p$ enters
only through two elementary inequalities on $\N$, $4v\leq p^{v}+3$ and the Legendre margin
$2v_p(n!)\leq n-1$; both fail at $p=2$.

\begin{theorem}[The $p$-adic binomial theorem at an arbitrary exponent]
  \label{thm:padic-binomial}
  \uses{def:padicexplog}
  \lean{LWX.PadicExpLog.hasSum_choose_mul_pow, LWX.PadicExpLog.chooseCoeff,
        LWX.PadicExpLog.norm_chooseCoeff_le,
        LWX.PadicExpLog.eq_zero_of_forall_tsum_natCast_eq_zero}
  \leanok
  For $e, x \in K$ with $\lVert x\rVert\leq c$, $\lVert e\rVert\lVert x\rVert\leq c$ and
  $c^{2}<\lVert p\rVert$,
  \[
    \sum_{r\geq 0}\binom{e}{r}x^{r} \;=\; \exp\bigl(e\log(1+x)\bigr).
  \]
\end{theorem}

\begin{proof}
  \leanok
  \uses{def:padicexplog}
  Fix $x$ and regard both sides as power series in the exponent $e$, with geometrically decaying
  coefficients.  They agree at every $e = n \in \N$, by the finite binomial theorem on the left
  and $\exp(n\log u) = u^{n}$ on the right.  A power series vanishing at all natural numbers
  vanishes identically, by Strassman's theorem applied along $p^{k}\to 0$.
\end{proof}

The exponent has to be arbitrary: in the application it is either the normalised logarithm
$\ell\langle x\rangle$ of a $1$-unit, which need not lie in $\Zp$ when $K/\Qp$ is ramified, or
the halo exponent $s = \log(1+T_0)/p$, whose norm $p\lVert T_0\rVert$ exceeds $1$.  So the
density-of-$\N$ proof used in Chapter~\ref{chap:jacobs} does not apply, and the identity theorem
above replaces it.

\begin{lemma}[The binomial power identity]
  \label{lem:pow-sub-one}
  \uses{thm:padic-binomial}
  \lean{LWX.oneAddPow_pow_mul, LWX.norm_pow_prime_pow_sub_one_le,
        LWX.tendsto_pow_prime_pow_sub_one, LWX.exists_sq_norm_pow_prime_pow_sub_one_lt}
  \leanok
  $(1+T)^{p^{h}u} = \bigl((1+T)^{p^{h}}\bigr)^{u}$ for $u\in\Zp$; and
  $(1+T)^{p^{h}} - 1 \to 0$ as $h\to\infty$.  Hence every halo point reaches the joint
  $\exp/\log$ disc at some analyticity level.
\end{lemma}

\begin{proof}
  \leanok
  Both sides of the identity are continuous in $u$ and agree on $\N$, where they are the ordinary
  binomial theorem, and $\N$ is dense in $\Zp$.  For the decay, iterate the one-step Kummer bound
  $\lVert (1+T)^{p}-1\rVert \leq \max(\lVert p\rVert\lVert T\rVert, \lVert T\rVert^{p})$, whose
  contraction factor is $\max(\lVert p\rVert, \lVert T\rVert^{p-1})<1$.
\end{proof}

The second half is the quantitative form of \cite[\S2.7]{LWX}'s statement that the universal
character is $m$-locally analytic for $m$ large enough, and it is what lets
Theorem~\ref{thm:lwx-prop217-h} cover the whole annulus.

\section{The halo weight}
\label{sec:haloweight}

\begin{definition}
  \label{def:haloweight}
  \uses{def:univchar, thm:padic-binomial, def:analyticweight}
  \lean{LWX.haloUnits, LWX.haloChar, LWX.M1K, LWX.haloRho, LWX.haloCol, LWX.haloExpansion,
        LWX.haloWeight, LWX.evalAt_autFactor_haloWeight}
  \leanok
  The specialisation of the universal character at $T_0$ is
  $\kappa_{T_0}(x) = \omega(\bar x)\,(1+T_0)^{\ell\langle x\rangle}
  = \omega(\bar x)\exp\bigl(s\log\langle x\rangle\bigr)$, with halo exponent
  $s = \log(1+T_0)/p$.  On the sub-annulus $p^{-1}<\lVert T_0\rVert$, $\lVert T_0\rVert^{2}<p^{-1}$
  it is an analytic weight in the sense of Definition~\ref{def:analyticweight}: the group is
  $\psi(\Zp^{\times})(1+p\calO_K)$, the level is the image of $M_1$, and the expansion of
  $\kappa_{T_0}(cz+d)$ on the closed unit ball is
  \[
    (cz+d)^{2}\,\kappa_{T_0}(d)\sum_{m\geq 0}\binom{s}{m}\Bigl(\frac{c}{d}\Bigr)^{m}z^{m},
    \qquad \text{decay } \rho = \lVert T_0\rVert\sqrt{p} .
  \]
\end{definition}

The factor $(cz+d)^{2}$ is the normalisation of Definition~\ref{def:jacobs-weight}, chosen so
that the automorphy factor $\kappa(cz+d)/(cz+d)^{2}$ of Definition~\ref{def:kappaslash} is
literally \cite[(2.3.2)]{LWX}'s $\chi(cz+d)$.  The expansion is
Theorem~\ref{thm:padic-binomial}, and the constraint $\lVert T_0\rVert^{2}<p^{-1}$, that is
$v(T_0)>\tfrac12$, is exactly the condition for its convergence.  \cite[\S2.7]{LWX} record only
the cruder radius and remark that ``the radius here is not optimal''.

\begin{definition}[The halo weight at analyticity level $h$]
  \label{def:haloweighth}
  \uses{def:haloweight, lem:pow-sub-one}
  \lean{LWX.Mh, LWX.M1Kh, LWX.TH, LWX.haloExponentH, LWX.haloCharH, LWX.haloRhoH,
        LWX.haloWeightH, LWX.specialize_univChar_eq_padicExp,
        LWX.evalAt_autFactor_haloWeightH}
  \leanok
  With $h = m-1$, $T'_h = (1+T_0)^{p^{h}}-1$ and $s_h = \log(1+T'_h)/p^{h+1}$, the character on
  the disc of radius $p^{-(h+1)}$ about a $\Zp$-unit $a$ is
  $x\mapsto [a](T_0)\exp\bigl(s_h\log(x/a)\bigr)$, an analytic weight at the level
  $M_1^{(h)}$ of matrices with $p^{h+1}\mid c$, of radius
  $\max(p^{-(h+1)}, \lVert T'_h\rVert)\sqrt p$.
\end{definition}

That the two descriptions agree, that is
$(1+T_0)^{\log u/p} = (1+T'_h)^{\log u/p^{h+1}} = \exp(s_h\log u)$ on $1+p^{h+1}\Zp$, is
Lemma~\ref{lem:pow-sub-one}.  This is \cite[\S2.1]{LWX}: ``$\chi(a\cdot x) =
\chi(a)\cdot\chi(\exp(p^{m}))^{(\log x)/p^{m}}$''.

\section{The integral model $S^{D}_{\mathrm{int}}$}
\label{sec:intmodel}

\cite[\S2.7]{LWX} define the space of integral $p$-adic automorphic forms as the
$\Iw_q$-equivariant functions valued in $\Ind^{\Iw_q}_{B(\Zp)}[-]$, identified by
\cite[(2.3.1)]{LWX} with $C(\Zp;\Lambda)$ carrying the action
\[
  (h\Vert_\delta)(z) \;=\; [\,cz+d\,]\cdot h\Bigl(\frac{az+b}{cz+d}\Bigr)
  \tag{2.3.2}
\]
of the monoid $M_1$.  The repository's Hecke layer (Chapter~\ref{chap:qmf}) is stated over an
arbitrary semiring, so it can be instantiated at $R = \Lhalo$ and reused rather than re-proved.

\begin{definition}
  \label{def:intforms}
  \uses{def:univchar, def:localmat, def:forms}
  \lean{LWX.CFun, LWX.cfunSlash, LWX.cfunSlashAction, LWX.IntForms, LWX.intAutSlashAction}
  \leanok
  $C(\Zp,\Lhalo)$ with the action (2.3.2) of $M_1$, and $S^{D}_{\mathrm{int}}$ the corresponding
  space of automorphic functions at a tame level $U$ --- the same construction as
  Definition~\ref{def:forms}, at the ring $\Lhalo$ instead of a field.
\end{definition}

\begin{proposition}[{\cite[Proposition 3.4]{LWX}}]
  \label{prop:prop34}
  \uses{def:intforms, def:entry}
  \lean{LWX.mahlerON, LWX.mahlerCoeffs, LWX.cfunSlash_mahlerON, LWX.entry_eq_fwdDiff,
        LWX.mahlerEmbed}
  \leanok
  The Mahler basis $\binom{z}{n}$ realises $C(\Zp,\Lhalo)$ as $c(\N,\Lhalo)$, and in these
  coordinates the matrix of $\Vert_\delta$ is exactly the entry stream $P_{m,n}(\delta)$ of
  Definition~\ref{def:entry}.  The $T$-rescaled submodule
  $\widehat{\bigoplus}_n T^{n}\Lhalo\binom{z}{n}$ of \cite[(5.4.1)]{LWX} is realised by an
  explicit embedding.
\end{proposition}

\begin{proof}
  \leanok
  \uses{def:entry, def:univchar}
  Mahler's theorem applies because $\Lhalo$ is an ultrametric $\Zp$-Banach algebra.  The Mahler
  coefficients of a function are its iterated forward differences at $0$, so the coefficient of
  $\binom{z}{n}$ in $\binom{z}{n'}\Vert_\delta$ is the $m$-th difference at $0$ of
  $\binom{f(z)}{n'}[cz+d]$, which is Definition~\ref{def:entry} verbatim; the $T$-expansion of
  $[cz+d]$ is the one performed there.
\end{proof}

So the definition of Chapter~\ref{chap:lwx-halo}, which was taken as a definition there, is a
theorem here.  Note that the rescaling of \cite[(5.4.1)]{LWX} lives in the estimates and not in
the model coordinates: the action (2.3.2) does not preserve $c_0$-coordinates in the rescaled
basis.

\begin{theorem}[{\cite[(2.11.1)]{LWX} and \cite[Proposition 3.1]{LWX}}]
  \label{thm:prop31}
  \uses{def:intforms, def:updatum, thm:evalatreps, def:heckeoperator}
  \lean{LWX.intEvalAtReps, LWX.intEvalAtReps_comm,
        LWX.bijective_intEvalAtReps_of_stabilizer_eq_bot, LWX.certM1, LWX.UpDatum.ofCerts,
        LWX.isUpShape_certM1, LWX.intHeckeOperator, LWX.intEvalAtReps_intHeckeOperator}
  \leanok
  At a family of class representatives with trivial stabilisers, evaluation at the
  representatives is a bijection $S^{D}_{\mathrm{int}} \to c(\iota\times\N,\Lhalo)$; and for the
  genuine integral Hecke operator $[U\eta U]$ the diagram
  \[
    \mathrm{ev}\circ[U\eta U] \;=\; D.\mathrm{op}\circ\mathrm{ev}
  \]
  commutes, where $D$ is the $U_p$-datum assembled from right-coset certificates of $U\eta U$.
  Its local matrices have $U_p$-shape because they lie in
  $\Iw_q\begin{pmatrix} p&0\\0&1\end{pmatrix}\Iw_q \subseteq
  \bigl(\begin{smallmatrix} p\Zp&\Zp\\ q\Zp&\Zp^{\times}\end{smallmatrix}\bigr)$.
\end{theorem}

\begin{proof}
  \leanok
  \uses{thm:evalatreps, def:heckeoperator, def:updatum}
  \cite[Proposition 3.1]{LWX} writes $\gamma_i v_j^{-1} = \delta_{i,j}^{-1}\gamma_{\lambda_{i,j}}
  u_{i,j}$ and finds $(U_p\varphi)(\gamma_i) = \sum_j \varphi(\gamma_{\lambda_{i,j}})
  \Vert_{u_{i,j,p}v_j}$.  That display is already proved in the ring-generic Hecke layer
  (Theorem~\ref{thm:evalatreps} and its companion for $[U\eta U]$), so what is left is to check
  that any operator satisfying it intertwines evaluation with $D.\mathrm{op}$, which is a
  computation with the block decomposition, and to verify the shape of
  $\delta_{i,j,p} = u_{i,j,p}v_j$ from the Iwahori double coset.
\end{proof}

\section{Amice's theorem and the Colmez basis}
\label{sec:amice}

\cite[\S2.16]{LWX} quote \cite[Th\'eor\`eme 1.29]{Colmez}: the functions
$\lfloor n/(q^{-1}p^{m})\rfloor!\binom{z}{n}$ form an orthonormal basis of the space
$\mathcal{O}B_{qp^{-m}}$ of functions analytic on the discs $a+q^{-1}p^{m}\Zp$.  This is Amice's
theorem \cite{Amice}.  We prove it at every level, in the form the seam needs.

\begin{proposition}[Level $m=1$]
  \label{prop:colmez-m1}
  \uses{def:tate-algebra, def:cspace}
  \lean{LWX.colmezToMonomial, LWX.norm_colmezToMonomial, LWX.colmezEquiv,
        LWX.monomialToMahler, LWX.diagFactorial, LWX.monomialToMahler_eq,
        LWX.mahlerCoeffPow, LWX.pow_eq_sum_choose_mul_mahlerCoeffPow,
        LWX.fwdDiff_iter_evalAt_natCast}
  \leanok
  At $m=1$ (so $p$ odd, $q=p$) the disc is the closed unit disc and the space is the Tate algebra
  $K\langle z\rangle = c(\N,K)$; the Colmez basis is
  $n!\binom{z}{n} = z(z-1)\cdots(z-n+1)$, a monic integral polynomial.  Change of coordinates
  from the Colmez basis to the monomial basis is an isometry with dense range, hence a continuous
  linear equivalence, and
  \[
    \mathrm{monomialToMahler} \;=\; \mathrm{diag}(m!)\circ\mathrm{colmezEquiv}^{-1} .
  \]
\end{proposition}

The last display is the ``conjugated by an infinite diagonal matrix'' of
\cite[proof of Proposition 2.17]{LWX} at $m=1$.  The isometry is proved by a largest-index
argument; the same argument at every level is Theorem~\ref{thm:unitriangular}.

\begin{lemma}[{\cite[Lemme 1.4.9]{Colmez}}]
  \label{lem:colmez-valuation}
  \lean{LWX.colmezPoly, LWX.discPoly, LWX.discFactor, LWX.discConst, LWX.discPoly_eq,
        LWX.norm_discConst, LWX.norm_coeff_discPoly_le, LWX.norm_coeff_discPoly_le_inv_of_lt,
        LWX.norm_coeff_discPoly_diag}
  \leanok
  Let $g_n = \lfloor n/p^{h}\rfloor!\binom{z}{n}$ and let $a < p^{h}$ be a disc centre.  The
  polynomial $g_{n,a}(w) = g_n(a+p^{h}w)$ has $\Zp$-coefficients, and factors as
  $g_{n,a} = c_{n,a}f_{n,a}$ where $f_{n,a}$ reduces mod $p$ to a monic polynomial of degree
  $\lvert K_{n,a}\rvert$, $K_{n,a} = \{k<n : k\equiv a \pmod{p^{h}}\}$, and
  $v(c_{n,a}) = \sum_{\ell=1}^{h}[\,a \bmod p^{\ell} < n\bmod p^{\ell}\,] \geq 0$.
\end{lemma}

\begin{proof}
  \leanok
  Each linear factor $a-k+p^{h}w$ is $p^{h}(w-(k-a)/p^{h})$ when $p^{h}\mid k-a$, and otherwise
  $(a-k)\bigl(1+p^{h}w/(a-k)\bigr)$; collecting the two kinds of factor gives the factorisation,
  and the valuation of the scalar is the count of carries, which is the displayed sum.
\end{proof}

\cite{Colmez} centres the discs at $-j$ with $1\leq j\leq p^{h}$; we centre at the natural
representative $a = n\bmod p^{h}$, which is a translation of the unit disc and changes neither
integrality, nor vanishing mod $p$, nor degrees.

\begin{theorem}[Amice's theorem at level $h$]
  \label{thm:amice}
  \uses{lem:colmez-valuation, thm:unitriangular, def:cspace}
  \lean{LWX.discIdx, LWX.revIdx, LWX.colmezMatrix,
        LWX.isUnitriangularPerturbation_colmezMatrix, LWX.colmezToDisc,
        LWX.colmezDiscEquiv, LWX.diagFactorialH, LWX.discToMahler,
        LWX.discToMahler_apply_eq_fwdDiff}
  \leanok
  Model $\mathcal{O}B_{p^{-h}}$, the functions analytic on every disc $a+p^{h}\Zp$ with the
  maximum of the $p^{h}$ Gauss norms, by $c(\Z/p^{h}\times\N, K)$ --- the Taylor coefficients on
  each disc in the coordinate $z = a+p^{h}w$.  Then the map sending $e_n$ to the disc model of
  $g_n$ is an isometric equivalence $c(\N,K) \to c(\Z/p^{h}\times\N,K)$, and composing with
  $\mathrm{diag}(\lfloor n/p^{h}\rfloor!)$ turns Colmez coordinates into Mahler coefficients.
\end{theorem}

\begin{proof}
  \leanok
  \uses{lem:colmez-valuation, thm:unitriangular}
  Order the basis by $n = (m(n)+1)p^{h} - i(n)$ with $1\leq i(n)\leq p^{h}$, and the discs by
  their centres, as \cite{Colmez} does; in that order Lemma~\ref{lem:colmez-valuation} says the
  reduction of the basis-change matrix is block upper triangular with each diagonal block lower
  triangular with invertible diagonal entries.  Reindexing the columns by the involution
  reversing each block of $p^{h}$ consecutive integers turns the matrix into a perturbed
  unitriangular one, so Theorem~\ref{thm:unitriangular} applies.  The two reindexings cancel, so
  the resulting map sends $e_n$ to $g_n$ on the nose.  The final clause is the level-$h$ form of
  the statement that Mahler coefficients are forward differences of values at $\N$-points.
\end{proof}

\begin{definition}[The disc model of $S^{D,\dagger,m}$]
  \label{def:discforms}
  \uses{thm:amice, def:haloweighth, def:forms, thm:qmf-compact}
  \lean{LWX.discImage, LWX.discConj, LWX.discSlash, LWX.discSlashAction, LWX.DiscForms,
        LWX.discEvalAtReps, LWX.bijective_discEvalAtReps_of_stabilizer_eq_bot,
        LWX.discHeckeOperator, LWX.discHeckeBlockOp, LWX.isCompactoid_discHeckeBlockOp,
        LWX.discHeckeCharPowerSeries, LWX.evalT_discHeckeCharPowerSeries_eq_zero_iff}
  \leanok
  The action (2.3.2) does not preserve a single disc: it \emph{permutes} the discs.  Writing
  $t_a(w) = a+p^{h}w$ and $a' = \delta(a)\bmod p^{h}$, the matrix identity
  $\delta\, t_a = t_{a'}\,\delta'$ with $\delta' = t_{a'}^{-1}\delta t_a$ of level $h$ says that
  the action carries the Taylor expansion on $a'$ to the one on $a$ through the single-disc
  weight action of Definition~\ref{def:haloweighth}.  So the disc action is a block operator of
  $\kappa$-slashes permuted by $a\mapsto a'$, and it is a right action because the two data
  satisfy the cocycle rules.  Everything in Chapter~\ref{chap:qmf} then transcribes: the space
  $S^{D,\dagger,m}$, evaluation at representatives, the Hecke operator, its compactness, its
  Fredholm determinant, and the eigenvalue reading of the zeros.
\end{definition}

\section{Proposition 2.17: the seam}
\label{sec:seam}

\begin{theorem}[{\cite[Proposition 2.17]{LWX}} at $m=1$]
  \label{thm:lwx-prop217}
  \uses{prop:colmez-m1, def:haloweight, thm:prop31, lem:diag-intertwine,
        thm:charpowerseries-inv, def:heckecharpowerseries}
  \lean{LWX.specOp, LWX.charPowerSeries_specOp, LWX.monomialToMahler_comp_kappaSlash,
        LWX.monomialToMahlerBlock_comp_heckeBlockOp,
        LWX.specCharSeries_ofCerts_eq_heckeCharPowerSeries}
  \leanok
  On the sub-annulus $p^{-1}<\lVert T_0\rVert$, $\lVert T_0\rVert^{2}<p^{-1}$, the integral
  characteristic series evaluated at $T_0$ is the Fredholm determinant of $U_p$ on the
  overconvergent forms of the halo weight:
  \[
    \Char(P)(T_0) \;=\; \det\bigl(1 - X\,[U\eta U]\bigr)
    \quad\text{on } S^{D,\dagger,1}_{\kappa_{T_0}}(U).
  \]
\end{theorem}

\begin{proof}
  \leanok
  \uses{prop:colmez-m1, lem:diag-intertwine, thm:charpowerseries-inv, prop:prop34}
  Four steps.  (i) The specialised operator $P(T_0)$ on $c(\iota\times\N,K)$ has
  $c_n\bigl(P(T_0)\bigr) = c_n(P)(T_0)$, because specialisation is a continuous ring homomorphism
  (Definition~\ref{def:halo-specialize}) and the coefficients are limits of determinants of
  minors.  (ii) The \emph{seam identity}: monomial-to-Mahler coordinates intertwine the analytic
  action of a certificate matrix with $P(T_0)$ --- both are $f\mapsto \chi(cz+d)f(\text{M\"ob}\,z)$
  read at $\N$-points --- and hence blockwise for $[U\eta U]$.  (iii) By
  Proposition~\ref{prop:colmez-m1}, $\mathrm{monomialToMahler} = \mathrm{diag}(m!)\circ
  \mathrm{colmezEquiv}^{-1}$, so the Colmez-basis matrix $P'$ satisfies
  $\mathrm{diag}(m!)P' = P(T_0)\mathrm{diag}(m!)$, and Lemma~\ref{lem:diag-intertwine} gives
  $\det(1-XP') = \det(1-XP(T_0))$ --- with no inversion of the unbounded diagonal.  (iv) Basis
  independence (Theorem~\ref{thm:charpowerseries-inv}) gives
  $\det(1-XP') = \det(1-X[U\eta U])$.
\end{proof}

\begin{theorem}[{\cite[Proposition 2.17]{LWX}} at every level]
  \label{thm:lwx-prop217-h}
  \uses{thm:lwx-prop217, thm:amice, def:discforms, lem:pow-sub-one}
  \lean{LWX.discToMahler_comp_discSlash, LWX.discToMahlerBlock_comp_discHeckeBlockOp,
        LWX.specCharSeries_ofCerts_eq_discHeckeCharPowerSeries,
        LWX.exists_level_specCharSeries_eq_heckeCharPowerSeries}
  \leanok
  The same identity holds at every analyticity level $m = h+1$, with $S^{D,\dagger,m}$ in its
  disc model.  Consequently every point of the full boundary annulus
  $p^{-1}<\lVert T_0\rVert<1$ is covered, at a level supplied by Lemma~\ref{lem:pow-sub-one}.
\end{theorem}

\begin{proof}
  \leanok
  \uses{thm:amice, def:discforms, lem:diag-intertwine, thm:charpowerseries-inv}
  The proof of Theorem~\ref{thm:lwx-prop217} with the disc model in place of the Tate algebra:
  Theorem~\ref{thm:amice} supplies the Colmez equivalence and the diagonal
  $\mathrm{diag}(\lfloor n/p^{h}\rfloor!)$, the disc action is (2.3.2) pointwise so on Mahler
  coefficients it is the specialised entry stream, and the two determinant identities are as
  before.
\end{proof}

\begin{corollary}[{\cite[Lemma 4]{Bu04}; \cite[Definition 2.13]{LWX}}]
  \label{cor:level-independence}
  \uses{thm:lwx-prop217-h}
  \lean{LWX.discHeckeCharPowerSeries_eq_of_levels,
        LWX.discHeckeCharPowerSeries_zero_eq_heckeCharPowerSeries}
  \leanok
  The characteristic power series of $U_p$ on $S^{D,\dagger,m}$ is independent of $m$: both
  levels compute $\Char(P)(T_0)$.
\end{corollary}

This is a statement the $m=1$ seam could not even express, and it is what makes
\cite[Definition 2.13]{LWX}'s spectral curve well defined.

\section{The quaternionic case}

Chapter~\ref{chap:qmf}'s quaternionic instantiation works over a number field $F$ at a place $v$
with $K = F_v$; the integral model needs $\Zp$-entries.  At $F=\Q$, $v = (p)$ the two are
compared by mathlib's isomorphism, once one knows it is an isometry.

\begin{lemma}
  \label{lem:adic-completion-isometry}
  \lean{Padic.absNorm_asIdeal_primesEquiv_symm, Padic.norm_adicCompletionEquiv}
  \leanok
  Mathlib's continuous $\Q$-algebra isomorphism $\Qp \simeq \Q_v$ at the place $v = (p)$ is an
  isometry for the finite-place norm.
\end{lemma}

\begin{proof}
  \leanok
  The absolute norm of the ideal $v$ is $p$, so $\lVert p\rVert = p^{-1}$ on both sides; and the
  closed unit balls correspond, by mathlib's comparison of $\Zp$ with the ring of integers of the
  completion.
\end{proof}

This generalises to every prime the computation made at $p=3$ in Chapter~\ref{chap:jacobs}.

\begin{theorem}[The spectral reading at a halo point]
  \label{thm:lwx-spectral}
  \uses{thm:lwx-prop217, thm:lwx-prop217-h, lem:adic-completion-isometry, thm:serre12,
        def:rigidification, thm:qmf-eigenform-criterion}
  \lean{LWX.thetaInt, LWX.IntFormsQ, LWX.intHeckeUpQ,
        LWX.specCharSeries_ofCerts_eq_heckeCharPowerSeriesQ,
        LWX.evalT_specCharSeries_eq_zero_iff,
        LWX.specCharSeries_ofCerts_eq_discHeckeCharPowerSeriesQ,
        LWX.evalT_specCharSeries_eq_zero_iff_disc,
        LWX.exists_level_evalT_specCharSeries_eq_zero_iff}
  \leanok
  Let $D$ be a definite quaternion algebra over $\Q$ split at $p$, $U$ a tame level, and suppose
  the level is neat.  Then $S^{D}_{\mathrm{int}}$ is the integral model of
  Definition~\ref{def:intforms} for $D$, and for every halo point $p^{-1}<\lVert T_0\rVert<1$
  there is an analyticity level at which, for $a \neq 0$,
  \[
    \Char(P)(T_0)(a) = 0
    \quad\Longleftrightarrow\quad
    a^{-1} \text{ is a } U_p\text{-eigenvalue on } S^{D,\dagger,m}_{\kappa_{T_0}}(U).
  \]
\end{theorem}

\begin{proof}
  \leanok
  \uses{thm:lwx-prop217-h, lem:adic-completion-isometry, thm:qmf-eigenform-criterion}
  Lemma~\ref{lem:adic-completion-isometry} identifies the integral component map with the
  $K$-side one of Chapter~\ref{chap:qmf}, so $S^{D}_{\mathrm{int}}$ and Buzzard's
  $S^{D}_{\kappa}(U)$ are the two sides of Theorem~\ref{thm:lwx-prop217-h}.  Composing that
  identity with the eigenvalue criterion of Theorem~\ref{thm:qmf-eigenform-criterion} --- itself
  Serre's Theorem~\ref{thm:serre12} in the block model --- gives the displayed equivalence.  The
  level is supplied by Lemma~\ref{lem:pow-sub-one}.
\end{proof}

Theorem~\ref{thm:lwx-spectral} is the fibre of \cite[Definition 2.13]{LWX}'s spectral curve over
a halo point: from here on, every statement proved about $\Char(P)$ in
Chapter~\ref{chap:lwx-halo} or Chapter~\ref{chap:lwx-slopes} is a statement about the genuine
$U_p$ on a genuine space of overconvergent quaternionic forms.
